\documentclass{ximera}

\title{Velocity and Applications of Derivatives}
\author{MATH 425: Calculus I}

\begin{document}
\begin{abstract}
   Working with the peers in your group, solve the following problems. Make sure to show and justify all your work. Make sure everyone in the group understands the solution and participates. Be prepared to report your answers to the whole class. 
\end{abstract}
\maketitle


\begin{exercise}
    One of the longest stretches of straight (and flat) road in North America can be found on the Great Plains in the state of North Dakota on state highway 46, which lies just south of interstate highway I-94 and runs through the town of Gackle.  A car leaves time (at time $t=0$) and heads east on highway 46; its postion in mles from Gackle at time $t$ in minutes is given by the graph of the function in the figure.  Three important points are labelled on the graph; where the function looks linear, assume it is a straight line.
    \begin{center}
        \desmos{bwkl25a63u}{800}{600}
    \end{center}
    The graph of $y=s(t)$, the position of the car along highway 46, which tells its distance from Gackle, ND, at time $t$ minutes.
    \begin{enumerate}
        \item In everyday language, describe the behavior of the car over the provided time interval. It may help to break your description down into the behavior on four distinct intervals: $[0,57]$, $[57,68]$, $[68,104]$ and after 104.
        \item Find the slope of the line between the points $(0,0)$ and $(57, 63.8)$.  What are the units of this slope?  What does the slope represent?
        \item Find the velocity of the car on each interval: $[0,57]$, $[57,68]$, $[68,104]$ and after 104.
        \item What is the instantaneous rate of change of the car's position at the moment $t=80$?  Explain what this value means.
    \end{enumerate}
\end{exercise}

\begin{exercise}
    The position of a particle moving along a straight line is given by $s(t)=t^3-6t^2+9t$, where $t$ is measured in seconds and $s$ in meters.
    \begin{enumerate}
        \item Find the velocity of the particle at time $t$. (Use the limit definition to find the derivative.)
        \item What is the velocity at $t=2$ seconds?
        \item When is the particle momentarily at rest? 
        \item Find the acceleration at time $t$, then at $t=2$ seconds. (Use the limit definition to find the derivative.)
    \end{enumerate}
\end{exercise}

\begin{exercise}
    The temperature, $H$, in degrees Celsius, of a cup of coffee placed on the kitchen counter is given by $H=f(t)$, where $t$ is in minutes since the coffee was put on the counter.
    \begin{enumerate}
        \item Would you expect $f'(t)$ to be positive or negative?  why? 
        \item What are the units of $f'(5)$?
        \item Suppose that $|f'(5)|=0.9$ and $f(5)=61$.  Write a sentence which describes the information you know about the coffee in plain language.
    \end{enumerate}
\end{exercise}

Problems 1 and 3 adapted from Boelkins, et al., \textit{Active Calculus} and Problem 2 adpated from Stewart, et al., \textit{Calculus}.



\end{document}
